\documentclass[11pt,a4paper]{article}

\usepackage{geometry}
\geometry{left=18mm,right=18mm,top=19mm,bottom=22mm}
\usepackage{fontspec}
\usepackage{xeCJK}
\setmainfont{Noto Serif}
\setsansfont{Noto Sans}
\setmonofont{Noto Sans Mono}
\setCJKmainfont{Noto Serif CJK JP}
\setCJKsansfont{Noto Sans CJK JP}
\setCJKmonofont{Noto Sans Mono CJK JP}
\XeTeXlinebreaklocale "ja"
\XeTeXlinebreakskip=0pt plus 1pt

\usepackage{xcolor}
\definecolor{Main}{HTML}{222222}
\definecolor{Light}{HTML}{F3F3F3}
\definecolor{Mid}{HTML}{666666}
\definecolor{Rule}{HTML}{AAAAAA}
\usepackage{hyperref}
\hypersetup{colorlinks=true,linkcolor=Main,urlcolor=Main,citecolor=Main,pdfauthor={ChatGPT},pdftitle={12章 ソフトウェア品質 日本語まとめ}}
\usepackage{booktabs}
\usepackage{tabularx}
\usepackage{longtable}
\usepackage{array}
\usepackage{ragged2e}
\usepackage{enumitem}
\usepackage{titlesec}
\usepackage{fancyhdr}
\usepackage{lastpage}
\usepackage[most]{tcolorbox}
\usepackage{needspace}
\setlength{\headheight}{14pt}

\newcolumntype{Y}{>{\RaggedRight\arraybackslash}X}
\newcolumntype{L}[1]{>{\RaggedRight\arraybackslash}p{#1}}
\setlength{\parindent}{0pt}
\setlength{\parskip}{0.55em}
\renewcommand{\baselinestretch}{1.05}
\setlist[itemize]{leftmargin=1.4em,itemsep=0.15em,topsep=0.15em}
\setlist[enumerate]{leftmargin=1.6em,itemsep=0.15em,topsep=0.15em}

\titleformat{\section}{\Large\bfseries\sffamily}{\thesection}{0.7em}{}
\titleformat{\subsection}{\large\bfseries\sffamily}{\thesubsection}{0.7em}{}
\titleformat{\subsubsection}{\normalsize\bfseries\sffamily}{\thesubsubsection}{0.7em}{}
\titlespacing*{\section}{0pt}{1.2em}{0.35em}
\titlespacing*{\subsection}{0pt}{0.9em}{0.25em}
\titlespacing*{\subsubsection}{0pt}{0.7em}{0.15em}

\pagestyle{fancy}
\fancyhf{}
\lhead{\small 12 章 ソフトウェア品質}
\rhead{\small SWEBOK v4.0a 日本語まとめ}
\cfoot{\small \thepage/\pageref{LastPage}}
\renewcommand{\headrulewidth}{0.3pt}
\renewcommand{\footrulewidth}{0pt}

\newtcolorbox{claimbox}{colback=Light,colframe=Rule,boxrule=0.4pt,arc=1mm,left=1.0em,right=1.0em,top=0.7em,bottom=0.7em}
\newtcolorbox{todobox}{colback=white,colframe=Rule,boxrule=0.4pt,arc=1mm,left=1.0em,right=1.0em,top=0.7em,bottom=0.7em,title={TODO（図）},fonttitle=\bfseries\sffamily}
\newtcolorbox{notebox}{colback=Light,colframe=Rule,boxrule=0.4pt,arc=1mm,left=1.0em,right=1.0em,top=0.7em,bottom=0.7em,title={重要},fonttitle=\bfseries\sffamily}
\newcommand{\term}[2]{\textbf{#1}（#2）}

\begin{document}

{\Huge\bfseries\sffamily 12 章 ソフトウェア品質\par}
{\Large Software Quality の日本語まとめ\par}
{\normalsize SWEBOK Guide v4.0a Chapter 12 を初学者向けに再構成\par}
{\normalsize 作成日：2026 年 5 月 12 日\par}

\begin{claimbox}
\textbf{この資料の主張}\par
ソフトウェア品質は、単に「バグが少ないこと」ではない。品質とは、利害関係者のニーズを正しく表した要求に、ソフトウェア製品・作業成果物・開発プロセスがどの程度適合しているかを扱う概念である。したがって、品質はテスト工程だけで作り込むものではなく、要求、設計、実装、レビュー、構成管理、測定、是正処置、運用までを通じて管理する対象である。
\end{claimbox}

\begin{claimbox}
\textbf{この資料の読み方}\par
専門用語は原則として日本語で示し、必要に応じて英語を括弧内に併記する。英語名は、元資料やツール上の名称を確認したいときの手がかりとして残す。初学者が前提知識なしに読めるように、各節では「何のための概念か」「実務では何を確認すればよいか」を重視する。文体は、だである調に統一する。
\end{claimbox}

\tableofcontents
\newpage

\section*{全体概要：この章で学ぶこと}
\addcontentsline{toc}{section}{全体概要：この章で学ぶこと}
この章は、ソフトウェア品質を、製品そのもの、製品を作るプロセス、途中で作られる作業成果物、そして品質を確保する活動の観点から扱う。品質は、要求に合っているかだけでなく、その要求が利害関係者の本当のニーズを表しているかにも依存する。よって、品質を考えるときは、要求、測定、レビュー、テスト、構成管理、標準、倫理、リスクを切り離して扱うことはできない。

\textbf{到達目標}\par
この章を読み終えると、次のことができるようになる。
\begin{itemize}
  \item ソフトウェア品質、製品品質、プロセス品質、品質要求の違いを説明できる。
  \item 誤り、欠陥、故障の違いを説明できる。
  \item 品質コストを、予防、評価、不適合、手戻りの観点から説明できる。
  \item 品質マネジメント、品質保証、品質制御、検証と妥当性確認の役割を区別できる。
  \item 静的解析、動的解析、形式的解析、レビュー、監査、テストの使い分けを説明できる。
  \item 測定、欠陥分類、根本原因分析、是正・予防処置を品質改善に結びつけて説明できる。
\end{itemize}

\begin{todobox}
12 章の全体地図を入れる。\par
\textbf{画像生成用プロンプト}\par
白背景、モノクロ、A4 本文幅に収まる横長の学習用図解を作成する。中央に「ソフトウェア品質」を置き、周囲に「品質の基礎」「品質マネジメント」「品質保証」「品質ツール」を配置する。品質の基礎から「文化と倫理」「品質コスト」「標準・モデル・認証」「ディペンダビリティと完全性レベル」へ分岐し、品質マネジメントから「改善」「計画」「評価」「是正・予防」へ分岐し、品質保証から「準備」「プロセス保証」「製品保証」「V\&V とテスト」へ分岐する。日本語ラベル、細線、講義ノート風、余白を広めにする。
\end{todobox}

\section*{最初に必要な用語}
\addcontentsline{toc}{section}{最初に必要な用語}

\noindent\begin{tabularx}{\textwidth}{L{0.27\textwidth}Y}
\toprule
\textbf{用語} & \textbf{初学者向けの説明} \\
\midrule
ソフトウェア品質 & ソフトウェアが、明示された要求と暗黙のニーズを、指定された条件のもとでどの程度満たすかを表す概念である。 \\
品質要求 & 性能、信頼性、安全性、セキュリティ、使いやすさなど、機能の実現水準や制約を表す要求である。 \\
品質マネジメント & 組織やプロジェクトが品質を方向づけ、管理するための活動である。 \\
品質保証 & プロセスや成果物が適切な品質を生み出すという信頼を得るための活動である。 \\
品質制御 & 成果物を測定・評価し、欠陥や問題を見つける活動である。 \\
検証 & 作ったものが、前工程や仕様で課された条件を満たしているかを確認することである。 \\
妥当性確認 & 作ったものが、意図した利用目的と利用者ニーズに合っているかを確認することである。 \\
ディペンダビリティ & 信頼して使える性質のまとまりであり、可用性、信頼性、保守性、安全性、セキュリティなどを含む。 \\
\bottomrule
\end{tabularx}


\section*{導入：品質はなぜ複数の知識領域にまたがるのか}
\addcontentsline{toc}{section}{導入：品質はなぜ複数の知識領域にまたがるのか}
ソフトウェア品質という語は、複数の意味で使われる。第一に、ソフトウェア製品に望まれる性質を指す。第二に、ある製品がその性質をどの程度持っているか、すなわち\term{ソフトウェア製品品質}{software product quality}を指す。第三に、その品質を得るためのプロセス、ツール、技法、すなわち\term{ソフトウェアプロセス品質}{software process quality}を指す。

品質は「要求への適合」と表現されることがある。ただし、ここでいう要求は、単に文書に書かれた要求だけでは不十分である。書かれた要求が、利害関係者のニーズ、期待、利用状況を正しく表していなければ、要求に適合してもよい品質とは言えない。

\term{品質要求}{quality requirement}は、機能要求に対する制約である。たとえば「検索する」という機能要求に対して、「通常時 2 秒以内に結果を返す」「障害時にもサービスを継続する」「個人情報を権限なく閲覧できない」といった品質要求が付く。品質要求は第 1 章の用語ではサービス品質制約に近い。

アジャイル開発や DevOps は、短い反復、早いフィードバック、利用者と開発者の近接、テストや運用の自動化によって、プロセス品質と製品品質の向上をねらう。品質は最後に検査で付け足すものではなく、開発と運用の流れ全体に組み込むべきものである。

\begin{todobox}
「品質」「品質要求」「製品品質」「プロセス品質」の関係図を入れる。\par
\textbf{画像生成用プロンプト}\par
白背景、モノクロ、A4 本文幅の概念図を作成する。中央に「ソフトウェア品質」を置き、左に「製品品質：製品が要求とニーズを満たす度合い」、右に「プロセス品質：製品を作る過程の良さ」、下に「品質要求：機能への制約・品質水準」を配置する。矢印で、品質要求が製品品質の判定基準になり、プロセス品質が製品品質に影響することを示す。日本語ラベル、講義ノート風、細線。
\end{todobox}

\section{ソフトウェア品質の基礎}
品質を確保するには、関係者全員が「何を品質とみなすか」を共有する必要がある。品質の合意が曖昧なままでは、開発者は何を優先して設計・実装・レビューすべきか判断できない。

品質上の典型的な課題には、要求を明確に定義しにくいこと、顧客や利用者とのコミュニケーション不足、仕様からの逸脱、アーキテクチャや設計の誤り、コーディング誤り、プロセスや手順への不適合、不十分なレビューやテスト、文書の誤りがある。

品質を扱ううえで、次の三つの用語を区別することが重要である。

\noindent\begin{tabularx}{\textwidth}{@{}L{0.16\textwidth}L{0.18\textwidth}Y@{}}
\toprule
\textbf{用語} & \textbf{英語} & \textbf{意味} \\
\midrule
誤り & error & 人間の行為が間違った結果を生むこと。例として、要求を誤解して設計する、条件式を取り違えるなどがある。 \\
欠陥 & defect / fault & 作業成果物の中に埋め込まれた不完全さである。要求、設計、コード、テスト文書などが仕様や要求を満たさない状態を指す。 \\
故障 & failure & 実行時や利用時に、システムが必要な機能を果たせない、または指定された範囲内で動作できない外部から見える現象である。 \\
\bottomrule
\end{tabularx}


\begin{notebox}
誤りは人間の行為であり、欠陥は成果物の中に潜む問題であり、故障は実行時に外から観測される問題である。レビューは欠陥を実行前に見つける活動であり、テストは故障を観測して欠陥の存在を示す活動である。
\end{notebox}

\begin{todobox}
誤り・欠陥・故障の関係図を入れる。\par
\textbf{画像生成用プロンプト}\par
白背景、モノクロ、A4 本文幅の横長図を作成する。左から「人間の誤り」「作業成果物に欠陥が埋め込まれる」「実行時に故障として現れる」を箱で並べ、矢印で接続する。下に「レビューは欠陥を早く発見」「テストは故障を観測」と注記する。日本語ラベル、講義ノート風。
\end{todobox}

\subsection{ソフトウェア工学文化と倫理}
組織文化は品質に強く影響する。たとえば、受託開発、組込み開発、量販ソフトウェア、クラウドサービスでは、品質への期待、リリース方法、レビューの重さ、運用責任が異なる。したがって、品質活動は一つの型をすべての組織へ機械的に当てはめるものではない。

健全なソフトウェア工学文化では、費用、納期、品質のトレードオフを正直に扱う。品質上のリスクや制約を隠すことは、短期的には進捗がよく見えるかもしれないが、長期的には利用者、顧客、社会に損害を与える。倫理は品質の一部である。技術者は、品質に関する情報、状況、結果を正確に報告する責任を持つ。

\subsection{品質の価値と費用}
品質保証は費用がかかるため、軽視されることがある。しかし、品質活動を行わないことにも費用がかかる。むしろ、欠陥が後工程や納品後に見つかるほど、修正費用と社会的影響は大きくなる。

\term{ソフトウェア品質コスト}{Cost of Software Quality, CoSQ}は、品質のために支払う費用と、品質が不足したために支払う費用を合わせて考える枠組みである。


\noindent\begin{tabularx}{\textwidth}{L{0.24\textwidth}Y}
\toprule
\textbf{費用区分} & \textbf{説明} \\
\midrule
実装費用 & 計画、設計、開発など、通常の開発作業にかかる費用である。 \\
予防費用 & 欠陥を作り込まないための費用である。プロセス改善、ツール、教育、標準整備などが該当する。 \\
評価費用 & 欠陥を見つけるための費用である。レビュー、監査、テスト、検査などが該当する。 \\
不適合・手戻り費用 & 見つかった欠陥を直す費用である。納品前の内部故障費用と、納品後の外部故障費用に分けられる。 \\
\bottomrule
\end{tabularx}


\term{適合費用}{conformance cost}は、予防費用と評価費用の合計である。これは品質を作り込み、確認するために意図的に投資する費用である。\term{不適合費用}{nonconformance cost}は、欠陥や故障が見つかったあとに発生する費用である。納品後の故障は、修正費用だけでなく、顧客の生産性低下、データ損失、評判低下、環境や公共への影響も含み得る。

\begin{todobox}
ソフトウェア品質コストの構造図を入れる。\par
\textbf{画像生成用プロンプト}\par
白背景、モノクロ、A4 本文幅の階層図を作成する。最上位に「ソフトウェア品質コスト」。第 2 層に「適合費用」と「不適合費用」。適合費用から「予防費用」「評価費用」へ分岐し、不適合費用から「納品前の手戻り」「納品後の外部故障」へ分岐する。下部に「最適な品質コストは、指定された品質水準を最小総費用で達成する点」と注記する。日本語、講義ノート風。
\end{todobox}

\subsection{標準・モデル・認証}
標準、モデル、認証は、品質活動を個人の経験だけに依存させないために役立つ。たとえば、ISO/IEC/IEEE 12207 はソフトウェアライフサイクルプロセスを扱う。安全性が重要な領域では、ソフトウェア安全計画や信頼性に関する標準が使われることがある。

一方、PDCA、CMMI、COBIT、PMBOK、BABOK、TOGAF などは、それぞれ異なる観点からプロセスやマネジメントの実践を整理する。ISO 9001 は品質マネジメント、ISO 27001 は情報セキュリティ、ISO 20000 はサービスマネジメントに関係する。スクラムや SAFe の認定は、アジャイル型の進め方に関する知識を示すことがある。

\begin{notebox}
標準や認証は、品質を自動的に保証する魔法ではない。重要なのは、対象の業界、リスク、契約、利用者への影響に合う標準を選び、実際のプロセスと成果物に反映することである。
\end{notebox}

\subsection{ソフトウェアのディペンダビリティと完全性レベル}
安全重要システムでは、ソフトウェアの故障が人命、環境、設備、公共サービスに大きな損害を与える可能性がある。このようなシステムでは、品質活動は通常の業務アプリケーションよりも厳密に計画される。

安全重要ソフトウェアには、システムに直接組み込まれる\textbf{直接ソフトウェア}と、安全重要ソフトウェアを開発・試験するために使われる\textbf{間接ソフトウェア}がある。航空機の飛行制御ソフトウェアは直接ソフトウェアの例である。安全重要ソフトウェアの開発環境やテスト環境のツールは間接ソフトウェアの例である。

故障リスクを下げる代表的な方法は、次の三つである。
\begin{itemize}
  \item \textbf{回避}：欠陥を作り込まないようにする。
  \item \textbf{検出と除去}：レビュー、解析、テストで欠陥を早く見つけて取り除く。
  \item \textbf{損害限定}：故障が起きても被害が広がらないようにする。
\end{itemize}

\term{保証ケース}{assurance case}は、ある性質が満たされているという主張を、根拠、論理、証拠によって説明する成果物である。安全性やセキュリティのように説明責任が重い品質では、単に「テストした」だけでなく、「なぜ十分だと考えるのか」を構造化して示す必要がある。

\subsubsection{ディペンダビリティ}
\term{ディペンダビリティ}{dependability}は、システムを信頼して使えることに関わる品質特性の集合である。可用性、信頼性、保守性、支援性、安全性、セキュリティなどが含まれる。故障が重大な結果を生むシステムでは、基本機能に加えてディペンダビリティが中心的な品質要求になる。

\subsubsection{ソフトウェア完全性レベル}
\term{完全性レベル}{integrity level}は、複雑さ、重要性、リスク、安全性、セキュリティ、性能、信頼性などから、システムやソフトウェアの重要度を示す値である。完全性レベルが高いほど、レビュー、検証、妥当性確認、独立性、証拠の厳密さが強く求められる。

完全性レベルは固定ではない。設計、実装、手順、技術上の対策によって上がることも下がることもある。取得者、供給者、開発者、独立した保証機関の合意によって決まることが多い。

\begin{todobox}
ディペンダビリティと完全性レベルの関係図を入れる。\par
\textbf{画像生成用プロンプト}\par
白背景、モノクロ、A4 本文幅の概念図を作成する。左に「利用状況・リスク・安全性・セキュリティ・性能」を配置し、中央の「完全性レベル決定」へ矢印を向ける。右に「必要な V\&V 技法」「独立性」「レビュー厳密度」「証拠量」を配置し、完全性レベルから矢印を伸ばす。下に「可用性・信頼性・保守性・安全性・セキュリティ」を「ディペンダビリティ」としてまとめる。日本語、細線。
\end{todobox}

\section{ソフトウェア品質マネジメントプロセス}
\term{ソフトウェア品質マネジメント}{Software Quality Management, SQM}は、ソフトウェア品質に関して組織を方向づけ、管理する活動である。目的は、製品、サービス、品質マネジメントプロセスの実施が、組織やプロジェクトの品質目標を満たし、顧客満足を達成することである。

中心概念は\term{品質マネジメントシステム}{Quality Management System, QMS}である。QMS は、プロセス、プロセス所有者、プロセス要求、測定、フィードバック経路を定義する。QMS には、品質保証、検証と妥当性確認、レビュー、監査、ソフトウェア構成管理などが含まれる。

品質活動が有効になるには、経営層の支援が不可欠である。プロジェクトが QMS を理解し、必要な教育と資源を持ち、不適合が見つかったときに是正できる体制が必要である。

\begin{todobox}
QMS と SQM のフィードバックループを入れる。\par
\textbf{画像生成用プロンプト}\par
白背景、モノクロ、A4 本文幅の循環図を作成する。「品質方針」「プロセス定義」「実施」「測定」「評価」「是正・予防」「改善」を円環状に配置し、中央に「QMS」を置く。外側に「経営層の支援」「プロジェクト」「顧客満足」を配置する。日本語ラベル、講義ノート風。
\end{todobox}

\subsection{ソフトウェア品質改善}
\term{ソフトウェア品質改善}{Software Quality Improvement, SQI}は、品質を継続的に高める活動である。代表的な考え方には、ソフトウェアプロセス改善、シックスシグマ、リーン、カイゼン、PDCA サイクルがある。

品質改善では、欠陥を後から見つけて直すだけではなく、欠陥が発生しにくいプロセスへ変えることが重要である。改善の対象は、ライフサイクルプロセス、欠陥の分類・検出・除去・予防、個人の能力改善である。

\term{品質機能展開}{Quality Function Deployment, QFD}は、顧客の声を品質目標や設計要素へ展開するための方法である。\term{個人ソフトウェアプロセス}{Personal Software Process, PSP}は、個々の技術者が自分の作業を測定し、技能と知識を継続的に改善する考え方である。

\begin{todobox}
PDCA と欠陥予防の流れを入れる。\par
\textbf{画像生成用プロンプト}\par
白背景、モノクロ、A4 本文幅の円環図を作成する。「計画」「実行」「確認」「改善」を時計回りに配置し、各段階に「品質目標を定める」「レビュー・テストを実施」「測定と欠陥分類」「根本原因を除去」と短い説明を添える。中央に「継続的品質改善」と書く。日本語、講義ノート風。
\end{todobox}

\subsection{品質マネジメント計画}
品質マネジメント計画では、どの品質プロセスを使うか、どの標準やモデルを使うか、具体的な品質目標をどう定めるか、各目標達成にどの程度の工数を使うか、どの活動時点で品質活動を行うかを決める。

組織は、品質方針、品質目標、手順を含む QMS を確立しなければならない。QMS の実施責任と権限は明確である必要があり、品質活動がプロジェクト管理上の短期的圧力だけで歪められないようにすることも重要である。

計画や手順は、利用者であるプロジェクトメンバーが使える形で文書化する必要がある。形式的に存在するだけの手順は、品質改善に役立たない。どの活動で、何を確認し、どの証拠を残すかが分かることが重要である。

\subsection{品質マネジメントの評価}
QMS が存在しても、それが機能しているとは限らない。したがって、品質マネジメント活動は評価される必要がある。プロセス能力や成熟度を評価し、その結果を改善プログラムへつなげる。

評価では、品質測定値、欠陥分類、監査結果、レビュー結果、テスト結果、不適合の数、是正処置の進捗などを見る。評価結果は、単なる報告で終わらせず、実行可能な改善課題に変換することが重要である。

\subsubsection{ソフトウェア品質測定}
品質測定は意思決定のために行う。ソフトウェアが複雑になるほど、「動くかどうか」だけではなく、「どの程度よく動くか」「どの品質目標をどの程度満たすか」を測る必要がある。

品質測定は、次の判断に役立つ。
\begin{itemize}
  \item 製品が品質目標をどの程度満たしているか。
  \item 工数、費用、納期に関する管理上の判断。
  \item テストをいつ終了してリリースするか。
  \item プロセス改善の効果が出ているか。
\end{itemize}

代表的な測定値には、\term{誤り密度}{error density}、\term{欠陥密度}{defect density}、\term{故障率}{failure rate}がある。欠陥密度は、欠陥数をソフトウェア規模で割った値である。故障率や平均故障時間は、信頼性モデルを作り、今後の故障可能性やテスト終了判断に使われることがある。

測定値は、表や数字だけでなく、パレート図、実行図、散布図、管理図、傾向分析、信頼性モデルとして可視化すると、どこに改善資源を集中すべきか判断しやすい。

\begin{todobox}
品質測定ダッシュボードの概念図を入れる。\par
\textbf{画像生成用プロンプト}\par
白背景、モノクロ、A4 本文幅のダッシュボード風図を作成する。4 つの小パネルを配置し、「欠陥密度」「故障率」「レビュー指摘件数」「是正処置の未完了数」を表示する。右側に「意思決定：リリース可否、重点改善領域、テスト継続」を箇条書きで示す。日本語ラベル、講義ノート風。
\end{todobox}

\subsection{是正処置と予防処置の実施}
品質目標が達成されない場合には、\term{是正処置}{corrective action}と\term{予防処置}{preventive action}が必要である。是正処置は、発生した問題を直すための処置である。予防処置は、同じ問題が将来再発しないように原因を取り除く処置である。

このためには、プロセスやツールの問題を報告できる仕組み、独立した記録・監視の仕組み、関係者へ進捗を知らせる仕組みが必要である。

\subsubsection{欠陥特徴づけ}
欠陥特徴づけとは、発見した誤り、欠陥、故障を分類し、原因や傾向を理解する活動である。単に「欠陥が何件あった」と数えるだけでは、再発防止に十分ではない。どの種類の欠陥が、どの工程で、どの成果物に、どの原因で入り込んだかを分類する必要がある。

\term{根本原因分析}{Root Cause Analysis, RCA}は、欠陥や不適合の表面的な現象ではなく、再発を生む根本原因を探す活動である。レビューやテストで発見された欠陥は、修正するだけでなく、プロセス、教育、ツール、標準、設計方針の改善へつなげるべきである。

\begin{todobox}
欠陥分類と根本原因分析の流れを入れる。\par
\textbf{画像生成用プロンプト}\par
白背景、モノクロ、A4 本文幅のフローチャートを作成する。「欠陥発見」から始まり、「記録」「分類」「傾向分析」「根本原因分析」「是正処置」「予防処置」「プロセス改善」へ進む。下部に「同じ欠陥を次のプロジェクトで繰り返さない」と注記する。日本語、細線、講義ノート風。
\end{todobox}

\section{ソフトウェア品質保証プロセス}
\term{ソフトウェア品質保証}{Software Quality Assurance, SQA}は、ソフトウェアプロセスが適切であり、意図された用途に適した品質のソフトウェア製品を生み出せるという信頼を確立するための活動である。

SQA は単なるテストではない。テストは重要だが、品質保証は、品質目標、プロセス、作業成果物、レビュー、監査、構成管理、測定、報告、是正処置を含む。特に重要システムでは、SQA 機能は技術的、管理的、財務的な圧力から独立していることが望まれる。

\begin{todobox}
SQA が「テストだけではない」ことを示す図を入れる。\par
\textbf{画像生成用プロンプト}\par
白背景、モノクロ、A4 本文幅の放射状図を作成する。中央に「SQA」を置き、周囲に「品質計画」「レビュー」「監査」「V\&V」「テスト」「構成管理」「測定」「是正処置」「教育」を配置する。下に「テストは SQA の一部である」と注記する。日本語ラベル、講義ノート風。
\end{todobox}

\subsection{品質保証の準備}
\term{ソフトウェア品質計画}{software quality plan}または\term{ソフトウェア品質保証計画}{Software Quality Assurance Plan, SQAP}は、特定の製品について、要求、利用者ニーズ、費用、納期、リスクに見合う品質を確保するための活動とタスクを定義する。

SQAP には、品質目標、品質活動、費用、資源、目的、スケジュール、参照する文書・標準・慣行・規約、測定値、統計的技法、問題報告と是正処置、ツール、技法、方法論、物理媒体のセキュリティ、教育、SQA 報告、文書化などが含まれる。

SQAP は、供給者ソフトウェアの調達、市販ソフトウェアの導入、納品後サービス、受け入れ基準なども扱うことがある。また、ソフトウェア構成管理計画と矛盾してはならない。構成管理、テスト、監査は相互補完的な活動である。

\begin{todobox}
SQAP の構成要素図を入れる。\par
\textbf{画像生成用プロンプト}\par
白背景、モノクロ、A4 本文幅のブロック図を作成する。中央に「SQAP」を置く。周囲に「品質目標」「標準・規約」「レビュー・監査」「測定」「問題報告」「是正処置」「ツール」「教育」「COTS・供給者」「受け入れ基準」「SCM との整合」を配置する。日本語ラベル、講義ノート風。
\end{todobox}

\subsection{プロセス保証の実施}
プロセス保証は、開発や保守のプロセスが、必要な品質の成果物を生み出すように定義され、実施され、記録されているかを確認する活動である。プロセス品質は最終製品品質に直接影響する。

プロセス保証では、組織の方針、プロセス、手順、役割、責任が明確であるかを見る。ライフサイクル活動の各節目で、文書レビュー、コードレビュー、V\&V、テスト、監査などの品質保証活動を計画し、実行し、記録する。

\term{ソフトウェア構成管理}{Software Configuration Management, SCM}もプロセス保証に関わる。構成項目の識別、変更管理、変更状態の記録、要求への適合確認を通じて、作業成果物とソフトウェアの品質を守る。

\subsection{製品保証の実施}
製品保証では、作るべきソフトウェアの本当の目的を確認する。利害関係者要求、機能要求、品質要求を正しく把握し、品質目標を測定可能な形で定義する必要がある。

ISO/IEC 25010 の品質モデルでは、ソフトウェア製品品質に関する代表的な特性として、機能適合性、性能効率性、互換性、使用性、信頼性、セキュリティ、保守性、移植性が扱われる。これらは単独で最適化できるとは限らない。たとえば、データ暗号化でセキュリティを高めると、性能に影響する場合がある。

製品保証の対象は、最終製品だけではない。要求仕様、設計記述、ソースコード、テスト文書、テスト報告書などの\term{作業成果物}{work product}も対象である。中間成果物の欠陥を早く見つけることで、後工程の大きな手戻りを防ぐ。

\begin{todobox}
プロセス保証と製品保証の違いを示す図を入れる。\par
\textbf{画像生成用プロンプト}\par
白背景、モノクロ、A4 本文幅の左右比較図を作成する。左に「プロセス保証：やり方が適切か」を置き、例として「手順」「役割」「SCM」「記録」「節目レビュー」を並べる。右に「製品保証：成果物が適切か」を置き、例として「要求仕様」「設計」「コード」「テスト文書」「最終製品」を並べる。中央に「両方が品質に必要」と注記する。日本語、講義ノート風。
\end{todobox}

\subsection{V\&V とテスト}
\term{検証}{verification}は、「正しく作っているか」を確認する活動である。あるライフサイクル段階の出力が、その段階の開始時に課された条件を満たしているかを見る。

\term{妥当性確認}{validation}は、「正しいものを作っているか」を確認する活動である。ソフトウェアが意図された利用目的や利用者ニーズを満たすかを見る。

V\&V の目的は、ライフサイクル全体を通じて品質を作り込むことである。テストは V\&V に含まれる重要な活動である。しかし、多くの場合、テストだけでは「意図した用途に適している」という十分な信頼を確立できない。要求が正確、完全、一貫、テスト可能であるかを早期に評価する必要がある。

重要システムでは、\term{独立検証妥当性確認}{Independent Verification and Validation, IV\&V}が行われることがある。これは、開発組織から技術的、管理的、財務的に独立した組織が V\&V を行うことである。

\subsubsection{静的解析技法}
\term{静的解析}{static analysis}は、ソフトウェアを実行せずに、要求、インタフェース仕様、設計、モデル、ソースコードなどの内容や構造を調べる技法である。コードリーディング、ピアレビュー、制御フローの静的解析、静的解析ツールによる検査などが含まれる。

実行されないコードは、動的テストでは直接確認できない。したがって、静的解析は、到達不能コード、構造的な問題、規約違反、潜在的な脆弱性を発見するうえで重要である。

\subsubsection{動的解析技法}
\term{動的解析}{dynamic analysis}は、ソフトウェアを実行またはシミュレーションして、誤りや欠陥を探す技法である。多くのテスト技法は動的解析である。ブラックボックステストでは、入力を与え、その出力を分析して要求への適合を判断する。

シミュレーション、モデル解析、モデル検査も、対象を実行・探索する意味で動的解析に含まれることがある。

\subsubsection{形式的解析技法}
\term{形式的解析}{formal analysis}は、数学的に定義されたモデルや仕様を用いて、誤りや不整合を探す方法である。安全性やセキュリティのように、特定性質を強く保証したい場合に使われることが多い。

形式的解析は強力だが、専門知識とコストが必要である。したがって、全体に一律適用するのではなく、リスクの高い部分や完全性レベルの高い機能に重点的に使うのが現実的である。

\subsubsection{品質制御とテスト}
\term{ソフトウェア品質制御}{Software Quality Control, SQC}は、プロジェクト成果物の品質を測定、評価、報告する活動である。テストは代表的な品質制御活動であり、開発成果物が入力要求を満たしていることを確認する検証活動の一部である。

品質活動では、テストレベル、テスト技法、測定、ツール、完了基準を計画する必要がある。重要システムでは、開発チームから独立した評価や、目標環境での適格性確認が求められることがある。

\subsubsection{技術レビューと監査}
レビューは、作業成果物の欠陥を早期に見つけるために有効である。コードが書かれたあとに欠陥を直すより、要求や設計段階で見つける方が安い。\term{ピアレビュー}{peer review}は、作業成果物を同僚が確認し、欠陥を除去するための活動である。近年の開発では、プルリクエストを使ったコードレビューがその一例である。

レビューには、目的、独立性、技法、役割、対象によってさまざまな種類がある。


\noindent\begin{tabularx}{\textwidth}{L{0.25\textwidth}Y}
\toprule
\textbf{レビュー種類} & \textbf{説明} \\
\midrule
アドホックレビュー & 構造化されておらず、レビュー担当者ができるだけ多くの欠陥を見つける。 \\
チェックリストベースレビュー & チェックリストに基づいて系統的に問題を探す。 \\
シナリオベースレビュー & 特定の読み方やシナリオに沿って成果物を確認する。 \\
観点ベースレビュー & セキュリティ、保守、利用者など特定の技術観点から確認する。 \\
役割ベースレビュー & 利害関係者の役割を想定して確認する。 \\
\bottomrule
\end{tabularx}


監査は、レビューより形式性が高く、独立性を確保するために第三者が行うことが多い。技術レビューや監査は、プロジェクト計画に組み込み、承認され、実施され、記録されるべきである。

\begin{todobox}
静的解析・動的解析・形式的解析の比較図を入れる。\par
\textbf{画像生成用プロンプト}\par
白背景、モノクロ、A4 本文幅の比較表風図を作成する。列は「静的解析」「動的解析」「形式的解析」。行は「実行するか」「主な対象」「代表例」「強み」「注意点」。各セルに短い日本語説明を入れる。講義ノート風、細線、印刷で読みやすい。
\end{todobox}

\begin{todobox}
レビューと監査の厳密さのスペクトラム図を入れる。\par
\textbf{画像生成用プロンプト}\par
白背景、モノクロ、A4 本文幅の横長の尺度図を作成する。左端を「軽量」、右端を「形式的」とし、左から「アドホックレビュー」「チェックリストレビュー」「シナリオレビュー」「観点ベースレビュー」「正式監査」を配置する。下に「目的、独立性、記録量が右へ行くほど大きくなる」と注記する。日本語、講義ノート風。
\end{todobox}

\section{ソフトウェア品質ツール}
ツールは品質活動を支援する。単純なものでは、要求追跡行列、レビュー用チェックリスト、コードレビュー用チェックリストがある。自動化ツールでは、Git のようなバージョン管理、プルリクエストによるコードレビュー、CI/CD、コード品質評価、自動テスト、オンデマンド環境が品質に貢献する。

品質ツールは、静的解析ツールと動的解析ツールに大きく分けられる。静的解析ツールは、ソースコード、モデル、文書を入力し、実行せずに構文的・意味的な分析を行う。動的解析ツールは、実行時の振る舞いを観察する。

代表的な品質ツールカテゴリは次のとおりである。

\noindent\begin{tabularx}{\textwidth}{@{}L{0.28\textwidth}Y@{}}
\toprule
\textbf{ツールカテゴリ} & \textbf{役割} \\
\midrule
レビュー・検査支援ツール & 文書やコードのレビュー作業を割り当て、コメント、欠陥記録、承認状態を管理する。 \\
安全ハザード分析支援ツール & FMEA や FTA などの安全分析を支援する。 \\
問題追跡ツール & テストや運用で見つかった異常を記録し、分析、処置、解決状態を追跡する。 \\
品質データ分析ツール & 欠陥密度、欠陥注入・除去率、傾向、予測をグラフや表で可視化する。 \\
CI/CD・自動テストツール & ビルド、テスト、デプロイ、品質ゲートを自動化し、早いフィードバックを与える。 \\
\bottomrule
\end{tabularx}


\begin{todobox}
品質ツールの分類図を入れる。\par
\textbf{画像生成用プロンプト}\par
白背景、モノクロ、A4 本文幅の三層図を作成する。上段に「品質ツール」。中段に「静的解析」「動的解析」「管理・可視化」。下段に、静的解析は「コード解析」「レビュー支援」「FMEA/FTA」、動的解析は「自動テスト」「性能測定」「シミュレーション」、管理・可視化は「問題追跡」「欠陥分析」「ダッシュボード」「CI/CD」を配置する。日本語ラベル、講義ノート風。
\end{todobox}

\section{トピックと参考文献の対応}
この表は、SWEBOK が示す「トピック対参考資料の行列」を、学習用に簡略化したものである。詳しく学ぶときは、各トピックに対応する文献や標準を読む。

{\small

\noindent\begin{longtable}{L{0.45\textwidth}L{0.48\textwidth}}
\toprule
\textbf{トピック} & \textbf{主な参考資料} \\
\midrule
\endhead
1. ソフトウェア品質の基礎 & Laporte and April, Software Quality Assurance; Sommerville, Software Engineering; IEEE/ACM Software Engineering Code of Ethics \\
1.1 ソフトウェア工学文化と倫理 & Laporte and April; IEEE/ACM Code of Ethics \\
1.2 品質の価値と費用 & Laporte and April; Crosby; Humphrey \\
1.3 標準・モデル・認証 & ISO/IEC/IEEE 12207; ISO 9001; CMMI; COBIT; PMBOK; BABOK; TOGAF \\
1.4 ディペンダビリティと完全性レベル & Sommerville; IEEE 1012; IEC 60300; DO-178C; EN 50128 \\
2. ソフトウェア品質マネジメントプロセス & Laporte and April; ISO/IEC/IEEE 90003; ISO/IEC TS 33061 \\
2.1 ソフトウェア品質改善 & PDCA; Lean; Kaizen; QFD; PSP; SPI \\
2.2 品質マネジメント計画 & Laporte and April; IEEE 730 \\
2.3 品質マネジメントの評価 & ISO/IEC TS 33061; IEEE 730; measurement literature \\
2.3.1 ソフトウェア品質測定 & Sommerville; IEEE 982.1; reliability models; Engineering Foundations KA \\
2.4 是正・予防処置 & QMS references; RCA references \\
2.4.1 欠陥特徴づけ & Laporte and April; defect taxonomy; RCA \\
3. ソフトウェア品質保証プロセス & IEEE 730; Laporte and April \\
3.1 品質保証の準備 & IEEE 730; SQAP guidance \\
3.2 プロセス保証の実施 & ISO 9001; Software Engineering Process KA; SCM KA \\
3.3 製品保証の実施 & ISO/IEC 25010; IEEE 982.1; Wiegers \\
3.4 V\&V とテスト & IEEE 1012; Software Testing KA; Formal Methods references \\
3.4.1 静的解析技法 & Review and audit standards; static analysis references \\
3.4.2 動的解析技法 & Software Testing KA; model analysis references \\
3.4.3 形式的解析技法 & Sommerville; Software Engineering Models and Methods KA \\
3.4.4 品質制御とテスト & IEEE 730; Software Testing KA \\
3.4.5 技術レビューと監査 & ISO/IEC 20246; IEEE 15288.2; Wiegers, Peer Reviews in Software \\
4. ソフトウェア品質ツール & Laporte and April; Software Construction KA; Software Testing KA; Software Operations KA \\
\bottomrule
\end{longtable}}

\section{発展的な読み物}
SWEBOK 12 章で紹介される発展的な読み物を、日本語で読む目的に合わせて整理する。


\noindent\begin{tabularx}{\textwidth}{L{0.35\textwidth}Y}
\toprule
\textbf{文献・標準} & \textbf{読むと得られること} \\
\midrule
IEEE 730-2014 & ソフトウェア品質保証プロセスを開始、計画、制御、実行するための要求を学べる。 \\
IEEE 1012:2016 & システム、ソフトウェア、ハードウェアの検証と妥当性確認を、完全性レベルに応じて計画する方法を学べる。 \\
ISO/IEC 20246:2017 & 作業成果物レビューの一般的な枠組みを学べる。 \\
Nancy Leveson, Safeware & ソフトウェア安全性の重要性と、安全性を開発プロジェクトへ組み込む考え方を学べる。 \\
Gilb and Graham, Software Inspection & レビューと欠陥に測定や統計的サンプリングを適用する考え方を学べる。 \\
Wiegers, Peer Reviews in Software & ピアレビューの種類、形式性、効果、実務導入方法を学べる。 \\
\bottomrule
\end{tabularx}


\section{参考文献}
\begin{enumerate}
  \item C. Y. Laporte and A. April, \textit{Software Quality Assurance}, IEEE Press, 2018.
  \item P. B. Crosby, \textit{Quality Is Free}, McGraw-Hill, 1979.
  \item W. Humphrey, \textit{Managing the Software Process}, Addison-Wesley, 1989.
  \item ISO/IEC 25010:2011, \textit{Systems and software Quality Requirements and Evaluation (SQuaRE) -- System and software quality models}, 2011.
  \item IEEE CS/ACM Joint Task Force, \textit{Software Engineering Code of Ethics and Professional Practice}.
  \item IEEE 730:2014, \textit{Standard for Software Quality Assurance Processes}, 2014.
  \item I. Sommerville, \textit{Software Engineering}, 10th ed., Addison-Wesley, 2016.
  \item RTCA, \textit{DO-178C: Software Considerations in Airborne Systems and Equipment Certification}, 2012.
  \item ISO/IEC 15026-1:2019, \textit{Systems and Software Assurance -- Concepts and Vocabulary}, 2019.
  \item ISO 9001:2015, \textit{Quality Management Systems -- Requirements}, 2015.
  \item IEEE 1012:2016, \textit{Standard for System and Software Verification and Validation}, 2016.
  \item ISO/IEC 20246:2017, \textit{Software and systems engineering -- Work product reviews}, 2017.
  \item K. E. Wiegers, \textit{Software Requirements}, 3rd ed., Microsoft Press, 2013.
  \item ISO/IEC/IEEE 24765:2017, \textit{Systems and Software Engineering -- Vocabulary}, 2017.
  \item N. Leveson, \textit{Safeware: System Safety and Computers}, Addison-Wesley Professional, 1995.
  \item T. Gilb and D. Graham, \textit{Software Inspection}, Addison-Wesley Professional, 1993.
  \item K. E. Wiegers, \textit{Peer Reviews in Software: A Practical Guide}, Addison-Wesley Professional, 2001.
  \item BS EN 50128:2011+A2:2020, \textit{Railway Applications -- Software for Railway Control and Protection Systems}, 2020.
  \item K. Iberle, \textit{They don't care about quality}, STAR East, 2013.
  \item D. Wallace, L. M. Ippolito, and B. B. Cuthill, \textit{Reference Information for the Software Verification and Validation Process}, NIST, 1996.
  \item IEC 60300-1:2014, \textit{Dependability Management -- Part 1}, 2014.
  \item D. Leffingwell, \textit{SAFe 4.5 Reference Guide}, Addison-Wesley, 2018.
  \item ISO/IEC TS 33061:2021, \textit{Process Assessment Model for Software Life Cycle Processes}, 2021.
  \item IEEE 15288.2:2014, \textit{Technical Reviews and Audits on Defense Programs}, 2014.
  \item PMI, \textit{A Guide to the Project Management Body of Knowledge}, 7th ed., 2021.
  \item ISO/IEC/IEEE 90003:2018, \textit{Guidelines for the application of ISO 9001:2015 to computer software}, 2018.
  \item COBIT, \textit{Control Objectives for Information Technology}, ISACA, 2019.
  \item IIBA, \textit{A Guide to the Business Analysis Body of Knowledge}, version 3, 2015.
  \item CMMI, \textit{Capability Maturity Model Integration}, ISACA, 2023.
  \item TOGAF, \textit{The Open Group Architecture Framework}, version 10, 2022.
  \item ISO/IEC 27001:2022, \textit{Information security management systems -- Requirements}, 2022.
  \item W. Humphrey, \textit{PSP: A Self-Improvement Process for Software Engineers}, Addison-Wesley Professional, 2005.
\end{enumerate}

\section{画像 TODO 一覧}
本文に配置した画像 TODO を再掲する。実際に図を作る場合は、各 TODO のプロンプトをそのまま画像生成に使うと、A4 資料に合う図を作りやすい。

\begin{enumerate}
  \item 12 章の全体地図。
  \item 品質・品質要求・製品品質・プロセス品質の関係図。
  \item 誤り・欠陥・故障の関係図。
  \item ソフトウェア品質コストの構造図。
  \item ディペンダビリティと完全性レベルの関係図。
  \item QMS と SQM のフィードバックループ。
  \item PDCA と欠陥予防の流れ。
  \item 品質測定ダッシュボード。
  \item 欠陥分類と根本原因分析の流れ。
  \item SQA がテストだけではないことを示す図。
  \item SQAP の構成要素図。
  \item プロセス保証と製品保証の違い。
  \item 静的解析・動的解析・形式的解析の比較図。
  \item レビューと監査の厳密さのスペクトラム図。
  \item 品質ツールの分類図。
\end{enumerate}

\section{章末確認問題}
\begin{enumerate}
  \item ソフトウェア品質、ソフトウェア製品品質、ソフトウェアプロセス品質の違いを説明せよ。
  \item 「要求に適合しているが、利用者の本当のニーズを満たしていない」場合、品質が高いと言えるか。理由を述べよ。
  \item 誤り、欠陥、故障の違いを、簡単な例を使って説明せよ。
  \item 品質コストにおける予防費用、評価費用、納品前手戻り費用、納品後外部故障費用の違いを説明せよ。
  \item QMS と SQM の関係を説明せよ。
  \item 品質改善で、欠陥を直すだけでなく根本原因分析が必要な理由を述べよ。
  \item SQA が単なるテストではない理由を説明せよ。
  \item 検証と妥当性確認の違いを説明せよ。
  \item 静的解析、動的解析、形式的解析の使い分けを説明せよ。
  \item レビューと監査の違いを説明せよ。
\end{enumerate}

\begin{claimbox}
\textbf{解答の方向性}\par
1 は「対象の違い」で説明する。2 は、要求がニーズを正しく表しているかに着目する。3 は「人間の誤り」「成果物中の欠陥」「実行時に観測される故障」の流れで説明する。4 は「品質へ投資する費用」と「品質不足で発生する費用」を区別する。7 は、SQA が品質計画、レビュー、監査、構成管理、測定、V\&V、是正処置を含むことを述べる。8 は、検証が「正しく作っているか」、妥当性確認が「正しいものを作っているか」であると説明できる。
\end{claimbox}

\end{document}
